\subsection{Optimal comparison constants and equality tangents}

The strategy of optimal comparison constants, local field changes,
and strict geodesic savings follows Gwynne and Miller
\cite[Section~1.5 and Sections~3--5]{GwynneMillerMetric}. The argument
below works with the conditional passage laws just constructed and
proves the required rigidity for those laws.

For a fixed joint law, define its deterministic optimal constants by
\begin{equation}\label{comp:optimal-constants}
 \begin{aligned}
 C&=\inf\{u>0:\text{upper comparison }u\text{ holds almost surely}\},\\
 c&=\sup\{u\ge0:\text{lower comparison }u\text{ holds almost surely}\}.
 \end{aligned}
\end{equation}
Both constants are deterministic functions of the joint law.
Countable rational port, cost, and confinement
tests make the events measurable.  Intersecting the probability-one
events for constants decreasing to $C$ and increasing to $c$, and
using arbitrary cost loss, shows that the optimal constants themselves
are valid whenever $C<\infty$.  The separate model initializations will
give $0\le c\le C<\infty$ and $C>0$ before rigidity is applied.

\begin{lemma}[Transport of comparison constants]
\label{comp:comparison-transport}
For common intrinsic kernels, almost-sure local comparisons on
corresponding buffered charts are unchanged by a fixed locally
equivalent field law.  Their optimal deterministic constants, defined
over the corresponding countable chart families, are unchanged as well.
They have the same numerical constants in the
indexed kernels of Equation~\eqref{comp:ann-index-definition}, when
both costs and reference lengths are expressed in the indicated
quantum units.
\end{lemma}

\begin{proof}
The joint density change is a function of the fields alone:
adjoining their same conditional passage kernel does not alter the
Radon--Nikodym derivative.  Thus a measurable local comparison event
of probability one keeps probability one.  Use a countable collection
of rational ports, confinements, and cost tolerances, and then continuum
joining, to obtain the simultaneous assertion.  Inverse absolute
continuity proves that a strict improvement under either law would
improve the other law, so the optimal constants coincide.
For a surface application, this last assertion uses the actual
chart-covering and quantum-unit identities stipulated above.  It does
not infer a comparison outside the tested charts from a single local
absolute-continuity statement.

For the lower constant, the local event is the intrinsic comparison.
An ambient lower bound $c>0$ implies it by
Lemma~\ref{comp:intrinsic-lower}; for $c=0$ it is automatic.
Conversely, divide a compact certificate into finitely many marked
pieces within a rational chart cover.  The local intrinsic bounds,
additive cost inequality, and ambient triangle inequality give the
same ambient lower bound for that certificate.  Hence the ambient
optimal definition and these measurable local tests are equivalent.

A fixed additive real-field change on the observation patch is locally
absolutely continuous under the stated Gaussian reference: extend it
smoothly across its buffer with support disjoint from the fixed pinning
circle.  The extension is a Cameron--Martin direction of that actual
Gaussian law.  In
Equation~\eqref{comp:ann-index-definition}, multiply the comparison
under $g+c_i$ by $v_i^{-p}$.  The reference Weyl identity gives
$v_i^{-p}D_{g+c_i}=D_g$.  The numerical comparison constant therefore
does not change.  Countably many fixed indices can be handled
simultaneously.  This argument transports an almost-sure assertion;
it does not identify individual passage samples before and after
the field change.
\end{proof}

We record explicitly the assumptions of the next implication.  An indexed
family has the \emph{annulus transfer property} if, for each fixed
$\vartheta>0$, sufficiently high probability of a buffered shell event
at at least $\vartheta N$ labels in $[N,2N]$ gives, with probability tending
to one, arbitrarily fine simultaneous good annuli at every point of a
fixed compact set.  The weight of a label $i$ is $w_i=\lambda^{-ip}$.
All reference regularity requirements below must be allowed in this
transfer.  In the application this is the absolute quantum-scale statement,
and it requires the common spatial passage kernels; it is not a
consequence of convergence of the contours alone.

\begin{proposition}[Amplification]
\label{comp:amplification}
Suppose the family has the annulus transfer property, the certificate
operations and continuum joining above, and finite optimal comparisons
$0\le c<C<\infty$.  Suppose the comparisons with the same numerical
constants hold at every indexed quantum scale and pass to further
extractions.  Then:
\begin{enumerate}[label=(\roman*)]
\item There are fixed shell collars and $b>0$ such that, for each
$\delta>0$ and all sufficiently large $N$, more than three quarters of
the labels $i\in[N,2N]$ have probability at least $b$ of a geodesic
traversal $P$ with
\[
 F_W(P(0),P(1))\ge(C-\delta)\operatorname{len}_D(P).
\]
Here $P$ crosses a fixed closed band, minimizes in its fixed open collar
$W$, and its entire range is contained in that closed band.
Its endpoints have fixed positive Euclidean separation.
The collars and $b$ do not depend on $\delta$.
\item For each $c_1\in(c,C)$ there are $b_1>0$, $\rho>0$ and fixed
rectangular collars such that more than three quarters of the labels
in every sufficiently late window have probability at least $b_1$ of
a rectangular certificate of cost
\[
 a<c_1D_{R^+}(x,y).
\]
Its endpoints lie on opposite faces of a rectangle $R$ of crossed width
at least $\rho$ and aspect ratio between $1/2$ and $2$, and its entire
range is contained in the closed rectangle $R$.  This closed envelope
is part of the retained witness.  Its slightly larger open collar is
used for certificate realizations, and $R^+$ is a further open
rectangular denominator collar.  One may impose deterministic positive lower
and finite upper bounds on the denominator, together with compact
internal metric moduli in a further collar.
\end{enumerate}
\end{proposition}

\begin{proof}
Fix the annulus threshold for favorable fraction $1/4$.  On a reference
regularity event of sufficiently high probability, all separated
crossings in the fixed shell cost at least $a_0>0$ in $D$, and surrounding
reference circuits have cost at most $A_0<\infty$.  These choices do not
involve $\delta$.  Choose $b$ smaller than one quarter of the allowable
failure probability.

If (i) fails, at least one quarter of the labels in arbitrarily late
windows have near-saturation probability less than $b$.  On the
complementary good event, every indicated geodesic traversal can be
followed at a saving of at least $\delta a_0w_i$ relative to the old
upper comparison.  Annulus transfer makes these events available
simultaneously at every point.  Stop a compact reference geodesic
successively when it exits good surrounding separators.  In each
completed interval retain the subarc from its last visit to the inner
testing circle until its first subsequent visit to the outer testing
circle.  This is minimizing in $W$, and has length at least $a_0w_i$.
The retained intervals are disjoint.

The reference separator chain gives
\[
 \operatorname{len}_D(P)\le A_0\sum_iw_i+o(1).
\]
Replacing the retained subarcs and following the gaps at upper factor
$C$ therefore gives total cost at most
\[
 C\operatorname{len}_D(P)-\delta a_0\sum_iw_i+o(1).
\]
At each fixed finite family of separators choose the total realization
and splicing error arbitrarily small.  Then shrink the separators.
This improves the upper factor to any number strictly above
$C-\delta a_0/A_0$, contradicting optimality.

For (ii), first choose a fixed small square size in the middle of the
shell.  On every hole-to-outer crossing, mark a point in this middle
region, follow the path until it first exits the square centered there,
and choose the exit face.  Start the marked piece at its last visit to
the parallel midline.  It lies in the resulting half-square and crosses
opposite faces; its width and its collar are fixed positive fractions
of the shell radius.  A high-probability reference event bounds the
$D$-distance of every pair of such opposite-face endpoints below by
$a_0w_i$: these pairs have uniformly positive Euclidean separation on
a compact set.

If shortcut probability were small on at least one quarter of the
labels, annulus transfer would make every marked rectangular piece
cost at least $c_1D_{R^+}(x,y)$.  For $c>0$, Lemma~\ref{comp:intrinsic-lower}
gives the old intrinsic lower bound on the same piece and on every gap;
for $c=0$ this assertion is automatic.  Thus its extra cost is at least
$(c_1-c)a_0w_i$.  The positive extra charge bounds the number of marked
pieces at a fixed finite annulus family even when $c=0$, by simultaneous
marked extraction and the additive cost inequality.  Reference
separator chaining now improves the lower comparison to any number
strictly below $c+(c_1-c)a_0/A_0$, a contradiction.

The marked half-square itself is retained as the closed envelope;
thus the shortcut event throughout this contradiction is the stronger
event in (ii).  Use an intermediate threshold between $c$ and $c_1$
before selecting rational outward collar widths and gridded parameter
bins.  Keep the actual envelope and its opposite-face endpoints inside
each bin, rather than replacing their range by the enlarged collar.
Strict slack preserves the inequality.  Finally, compactness of the family of reference
rectangles, continuity of $D$, and internal box connection bounds give
a reference regularity event of probability greater than $1-b_1/2$
with the additional bounds in (ii).  Intersecting it with each witness
event retains probability at least $b_1/2$.
\end{proof}

\begin{proposition}[Equality tangent]
\label{comp:equality-tangent}
In addition to Proposition~\ref{comp:amplification}, suppose the
common-kernel family is closed under joint extraction of its fixed
quantum offsets and its full marked certificate records, as in
Definition~\ref{common:admissible}(iv).  Suppose its
reference internal metrics on fixed collars retain their actual joint
law in these extractions.  Then a further law has, separately with
positive probability, a saturated nonconstant geodesic segment at upper
factor $C$ and a rectangular shortcut of ratio at most $C-\delta_0$,
where $\delta_0>0$ is deterministic.  The shortcut retains its closed
rectangular envelope and its endpoints on opposite envelope faces,
as in Proposition~\ref{comp:amplification}(ii).
Under local field equivalence and
strong confluence, the saturation event can be written for deterministic
anchors $a_0\ne a_1$, rational $0<t_1<t_2<1$, and a fixed bounded open
$U$ as
\begin{equation}\label{comp:fixed-saturation}
 P([t_1,t_2])\Subset U,\qquad
 F_U(P(t_1),P(t_2))=C(t_2-t_1)D(a_0,a_1),
\end{equation}
where $P$ is the unique fixed-anchor geodesic, parametrized by fractional
length.  Only the public strong-confluence assertion that a nonempty
geodesic subarc contains a segment of a geodesic between points of a
deterministic countable dense set is used.
\end{proposition}

\begin{proof}
Choose $c<c_1<c_2<C$ and $\delta_k\downarrow0$.  The two sets of labels
in Proposition~\ref{comp:amplification} intersect in more than half of
every sufficiently late window.  Choose a common label tending to
infinity.  Retain the two witnesses separately, using a dummy witness
off each event.  Their simultaneous occurrence is unnecessary.
The extraction property in Definition~\ref{common:admissible}(iv)
realizes the further law by a
diagonal discrete array with these witnesses, all fixed offsets, and
the common reference experiments retained.  Its kernels are the newly
extracted reference kernels; they need not be limits of kernel values
in the preceding laws.

The near-saturation paths have their entire ranges in a fixed closed
band $K$ and minimize in
its open collar $W$.  Retain $D_W|_{K\times K}$ as well as the paths.
On compact reference regularity sets, their lengths lie in a fixed
interval $[a,A]\subset(0,\infty)$ and their parametrizations have a
common Euclidean modulus.  Under a joint coupling, a subsequence
converges uniformly to a nonconstant $D_W$-geodesic $P$, since the
identities
$D_W(P_k(s),P_k(t))=|s-t|\operatorname{len}_D(P_k)$ pass to the
limit.  Exhausting these regularity sets loses arbitrarily little
probability.

If a limiting competitor in $W$ were strictly cheaper than
$C\operatorname{len}_D(P)$, its compact range would lie in a member
of a fixed countable exhaustion of $W$.  Hausdorff attainability of the
full certificate records realizes it in the approximating laws, with
nearby endpoints and almost the same cost.  The common upper comparison
and the continuum joining property correct those endpoints at vanishing
cost, uniformly on the retained compact reference regularity set.
This contradicts the near-saturation inequalities.  Thus the limiting
path is saturated.  Any cheaper competitor for an interior subsegment
could be joined to the unchanged remaining pieces at factor $C$,
contradicting saturation of the full path; every interior subsegment
is saturated as well.

For the low witness retain its closed rectangular envelope $R_k$,
its endpoints, its certificate range, and its two outer collars.
The envelope centers and side lengths lie in a compact parameter set
with positive lower width; the possible orientations form a finite
set.  Uniform reference moduli and the slack $c_1<c_2$ allow a finite
gridded library for the denominator alone: choose a fixed open
rectangle $\Omega$ inside the original denominator collar, still
containing $R_k$ with uniform positive clearance, and retain the
internal metric on $\Omega$ and a further fixed collar.
Shrinking the denominator can only increase its intrinsic distance.
Extract the actual envelope parameters as well as the marked
certificate and endpoints.  Along $R_k\to R$, closed range inclusion
and the opposite-face endpoint conditions pass to the limit.
The limiting range therefore remains inside $R$ itself, regardless
of the open collar used to realize it.  Its ratio relative to
$D_\Omega$ is at most $c_2=C-\delta_0$.  Retaining this envelope
will allow the forcing construction to preserve the certificate on a
thin real-field patch still contained in the denominator domain.
If its lower comparison were $C>0$,
Lemma~\ref{comp:intrinsic-lower} would contradict this witness.

Transfer to a whole-plane field in a buffered patch using the same
local kernels.  Choose its actual Gaussian reference to be the
whole-plane GFF $G_R$ pinned on $\partial B_R(0)$, with this fixed
circle strictly outside the buffered patch.  All later candidate
supports are compactly contained in that patch, after discarding
finitely many coarse scales.  Thus every catalogue shift has zero
average on the pinning circle and belongs to the Cameron--Martin
space of this one reference law.  The annular Markov trials use
Dirichlet interiors conditioned on their exterior harmonic parts;
no change of density across the pinning circle is required.
Choose an interior segment shorter than its
$D$-distance to the patch boundary; it is an ambient geodesic because
any exiting competitor is longer.  Strong confluence places a
nonconstant part of it on a geodesic between two points of a fixed
countable dense set.  Fixed-pair uniqueness holds simultaneously for
these pairs.  Choose rational fractional times inside the common
segment and an open finite union of rational disks with closure inside
the saturation neighborhood.  This is a countable union of events
covering a positive-probability event; one choice has positive
probability and gives Equation~\eqref{comp:fixed-saturation}.
\end{proof}

\subsection{Local forcing and rigidity}

Assume $c<C$ and fix the equality tangent from
Proposition~\ref{comp:equality-tangent}, with the deterministic anchors,
fractional times and open set $U$ in Equation~\eqref{comp:fixed-saturation}.
For each sample let $P$ be its fixed-anchor reference geodesic and set
\begin{equation}\label{comp:output-ratio}
 R=\frac{F_U(P(t_1),P(t_2))}{(t_2-t_1)D(a_0,a_1)},
\end{equation}
with $R=+\infty$ if either endpoint lies outside $U$.
The saturation event has positive probability and gives $R=C$.
We will construct output samples with $C-\epsilon_N\le R<C$, where
$\epsilon_N\downarrow0$, and show that their total probability stays
bounded below.  This is impossible for a random variable under one
fixed law.

The equality tangent also has a separate positive-probability
rectangular shortcut event.  It need not occur together with saturation.
We will sample a shortcut in a small patch, retain the old passage
observations outside a larger patch, and change the real field by a
function $f$ that is constant on the shortcut patch and zero on the
preserved exterior.  The change will force the new reference geodesic
through corridors near both ends of the shortcut, where short passage
joins make the shortcut available.  Write $D'=D_{h+f}$ for the output
reference metric, $P'$ for its fixed-anchor geodesic, and $F'$ for the
completed output passage rule.

There are two probabilistic requirements.  After restricting the source
event by a fixed likelihood bound, each output subprobability measure
must be dominated by the same fixed multiple of the original law.
Also, many modifications of one source sample must not produce too much
mass at the same output.  The next two lemmas isolate these requirements.
Their application will count small balls hit by the old saturated
geodesic on the source side and traversals of unchanged shells by the
new geodesic on the output side.

\begin{lemma}[Restricted Cameron--Martin transfer]
\label{comp:restricted-cm-transfer}
Let
\[
 \mu(dh,d\eta,dZ)=G(dh)\rho(d\eta)K(h,\eta;dZ)
\]
be a joint law on standard Borel spaces, with Gaussian real field
law $G$.  Fix a deterministic Cameron--Martin shift $f$.
Suppose two separated passage restrictions have a conditional
product kernel, the first depending on a patch where $f=0$ and
the second on a patch where $f=b$.  Retain the first from the
old sample and sample the second independently, conditional on
the fields, at the exact shifted field $h+b$.
Assume locality identifies their joint kernel with the
corresponding restrictions of $K(h+f,\eta;\cdot)$.

There is a stochastic completion of these restrictions into
an output $(h+f,\eta,Z')$ with unselected law
$G_f(dh')\rho(d\eta)K(h',\eta;dZ')$.
For every selection event $A$ in the enlarged source space
contained in
\[
 \exp\left\{W_f(h)+\tfrac12\|f\|_{\mathcal H}^2\right\}
 \le\Lambda,
\]
its output subprobability measure $\nu_f$ satisfies
$\nu_f\le\Lambda\mu$.
\end{lemma}

\begin{proof}
Use the target regular conditional law of the full passage
object given the two retained restrictions and the output
fields.  Before imposing $A$, the locality and product
hypotheses give exactly the claimed output law.  Its density
relative to $\mu$ is the Gaussian shift density
\[
 L_f(h')=\exp\{W_f(h')-\tfrac12\|f\|_{\mathcal H}^2\}.
\]
At $h'=h+f$ this is
$\exp\{W_f(h)+\tfrac12\|f\|_{\mathcal H}^2\}$.
Restriction to $A$ decreases the unselected output measure
and confines it to the set where this density is at most
$\Lambda$.  This proves the assertion, regardless of whether
$A$ is measurable from the output.
\end{proof}

\begin{lemma}[Weighted output multiplicity]
\label{comp:weighted-output-multiplicity}
Let $\mu$ be a probability measure.  For each $N$ let $I_N$
be finite, with deterministic weights $w_\beta>0$,
subprobability measures $\nu_\beta$, and measurable sets
$H_N,B_\beta$.  Suppose, with constants independent of $N$,
\[
 \sum_{\beta\in I_N}w_\beta\nu_\beta(1)\ge\alpha N,
 \qquad
 \nu_\beta\le\Lambda\mu,
 \qquad
 \nu_\beta(H_N\cap B_\beta)=\nu_\beta(1),
\]
and, $\mu$-almost surely on $H_N$,
\[
 \sum_{\beta\in I_N}w_\beta\mathbf 1_{B_\beta}\le K N.
\]
Then $\mu(H_N)\ge\alpha/(\Lambda K)$.
In particular these assumptions cannot hold for
$H_N=\{C-\epsilon_N\le R<C\}$ when $R$ is one fixed-law
random variable and $\epsilon_N\to0$.
\end{lemma}

\begin{proof}
Domination, support, and Tonelli's theorem give
\[
 \alpha N
 \le \sum_\beta w_\beta\nu_\beta(H_N\cap B_\beta)
 \le \Lambda\int\mathbf 1_{H_N}
              \sum_\beta w_\beta\mathbf 1_{B_\beta}\,d\mu
 \le \Lambda K N\mu(H_N).
\]
For the last assertion, $\mathbf 1_{H_N}\to0$ pointwise,
including on $\{R=C\}$, so bounded convergence yields a
contradiction.
\end{proof}

\begin{remark}[Source mass and observable output counts]
\label{comp:counting-checks}
Suppose every point of a positive-length middle arc is covered
by at least $cN$ good holes, and a hit ball of deterministic
weight $w_\beta$ accounts for at most $K_0w_\beta$ of arc
length.  Integrating along the arc gives an order-$N$ lower
bound for the weighted number of hits.  This coverage statement
must already include all likelihood clipping and other
selection restrictions.

For the output bound, require $B_\beta$ to be a measurable
condition on the output fixed-anchor geodesic, including a
traversal of new length at least $\kappa w_\beta$ in a
deterministic enlarged shell.  Bounded overlap at each spatial
radius, $O(N)$ radii, boundedly many aligned volume labels,
and a fixed finite recipe list imply
\[
 \sum_\beta w_\beta\mathbf 1_{B_\beta}
 \le K_1N\operatorname{len}_{D'}(P').
\]
A deterministic output length cutoff turns this into the
hypothesis of Lemma~\ref{comp:weighted-output-multiplicity}.
Such a cutoff follows from an old length cutoff and a uniform
bound on the shift amplitude.  The events $B_\beta$ cannot
refer to an old geodesic or an unobserved preimage.
\end{remark}

The remaining work is to construct outputs to which these two lemmas
apply.  A deterministic field modification will first force the new
geodesic through corridors near both ends of a preserved shortcut.
Uniform stability
as the modified patch shrinks will then place its fixed-anchor ratio
arbitrarily close to $C$.  Finally we will make the forcing geometry
available in the annuli supplied by Lemma~\ref{comp:ann-coverage}.

\begin{lemma}[Deterministic gate forcing]\label{comp:gate-forcing}
Let $D$ be a proper length metric on $\mathbb C$ inducing its usual
topology, let $K$ be a closed disk, and let $a_0,a_1\notin K$.
Let $\ell_j$ be the infimum of the $D$-lengths of paths from $a_j$ to
$\partial K$ before entering $K$, and choose exterior minimizers $x_j$.
Suppose
\[
 D(a_0,a_1)\ge \ell_0+\ell_1+g,\qquad g>0.
\]
Let $D_f=e^{\xi f}\cdot D$, where $f$ is bounded, continuous, and
supported in the interior of $K$.  Fix $B,e,d>0$, $C>0$, and
$0<\delta<C$, with
\[
 B+2e<g,\qquad 4Ce<\delta d.
\]
Assume the following deterministic conditions.
\begin{enumerate}[label=(\roman*)]
\item A competitor from $a_0$ to $a_1$ has $D_f$-length at most
      $\ell_0+\ell_1+B$.
\item There are effective open regions $O,T_0,T_1,N$, with incidence
      graph $O-T_0-N-T_1-O$.  The exterior of $K$ is contained in $O$,
      $f\ge0$ on $O$, and the outer gates
      $A_j=O\cap T_j$ admit $D$-connections to $x_j$ of length at most
      $e$.  Every curve between its first and last visits to
      $\partial K$ with $D_f$-length at most $B$ admits an itinerary
      in this graph.  Intervening visits to an isolated inner region
      are impossible.
\item There are $q_0,q_1\in\Omega$, with
      $N\Subset\Omega$, such that
      $d=D_b(q_0,q_1;\Omega)$, where $D_b=e^{\xi b}D$.
      Every point of the terminal gate $G_j=N\cap T_j$ can be joined
      to $q_j$ inside $\Omega$ at $D_b$-cost at most $e$, and
      $f\ge b$ on $N\setminus(T_0\cup T_1)$.
\item Every complete corridor passage of $D_f$-length at most $B$
      from its outer gate to its terminal gate, or in the reverse
      direction, contains a point
      $z$ with $D_f(z,q_j;W)\le e$.  Here $W$ is one fixed open
      neighborhood containing the two marker-to-$q_j$ connections
      and the preserved certificate in (v), all with compact
      confinement.  It need not contain the exterior region $O$
      or the anchor approaches.
\item A passage rule $F'$ has upper comparison $C$ against $D_f$
      inside $W$, admits concatenation with arbitrarily small open
      loss, and contains a preserved certificate from $q_0$ to $q_1$
      in $W$ with cost at most $(C-\delta)d$.
\end{enumerate}
Then every $D_f$-geodesic from $a_0$ to $a_1$ has a subsegment
with endpoints $z_0,z_1\in W$ such that
\[
 F'_W(z_0,z_1)<C D_f(z_0,z_1).
\]
The same assertion holds for approximate exterior minimizers after
subtracting their errors from $g$ and allocating them within $e$.
\end{lemma}

\begin{proof}
Write $P'$ for a new geodesic.  It must enter $K$, since otherwise
its length is an old path length, contradicting (i) and $B<g$.
Its portion between the first and last visits to $\partial K$ has
length at most $B$: its exterior prefix and suffix have costs at
least $\ell_0$ and $\ell_1$.

Gate contacts below may lie in the closures of the open overlaps.
The outer $D$-distance bounds extend to these contacts by continuity.
For terminal contacts, $\overline N\subset\Omega$, and the internal
metric $D_b(\cdot,\cdot;\Omega)$ is locally continuous: sufficiently
close points are joined by near-minimizing paths contained in a fixed
small ball of $\Omega$.  Thus its distance bound to $q_j$ also extends
from $G_j$ to $\overline G_j$.  Realizing these infimal bounds may add
an arbitrarily small error, which the two strict displayed budget
inequalities absorb.

The exterior gap also gives $D(x_0,x_1)\ge g$, by concatenating
the two minimizing exterior approaches with a connection between
$x_0$ and $x_1$.  The first corridor visited by $P'$ is $T_0$.
Indeed, if it were $T_1$, the preceding part lies in $O$, where
its old cost is at most its new cost.  Its cost is at most
$\ell_0+B$, since the final exterior approach costs at least
$\ell_1$.  Appending the gate correction and the old exterior
approach from $x_1$ would give an old path of length at most
$\ell_0+\ell_1+B+e$, a contradiction.  The last corridor is
similarly $T_1$.  A path remaining in $O$ is excluded in the same
way.  No excursion in $O$ can connect the two outer gates:
such an excursion, of cost at most $B$, would give
$D(x_0,x_1)\le B+2e<g$.

The finite incidence graph consequently contains an inner
$T_0-N-T_1$ passage along $P'$.  For precision, a finite subdivision
of the compact curve subordinate to the open region cover gives
a finite region itinerary; repetitions can be retained or erased.
There is no need to sum corrections over repeated gate visits.
The first full approach to $N$ uses $T_0$, and the last full exit
from $N$ uses $T_1$.  Condition (iv) supplies points $z_0,z_1$ on
these passages, in this order, with
$D_f(z_j,q_j;W)\le e$.

Between these two points there is a central passage from the last
$T_0$ contact to the first following $T_1$ contact.  This passage
lies in $N$, avoids the effective corridors in its interior, and
has multiplier at least $e^{\xi b}$.  Correcting its terminal
points inside $\Omega$ by (iii) shows that its new length is at
least $d-2e$.  In particular
\[
 D_f(z_0,z_1)=\operatorname{len}_{D_f}(P'[z_0,z_1])\ge d-2e.
\]
Upper comparison on the two short approach joins and the
preserved certificate now give
\[
 F'_W(z_0,z_1)
 \le (C-\delta)d+2Ce
 <C(d-2e)
 \le C D_f(z_0,z_1).
\]
The strict middle inequality leaves room for the arbitrarily
small concatenation losses.
\end{proof}

Figure~\ref{fig:forcing-corridors} displays the preservation patch,
two corridors, and transverse markers used in the construction below.
\begin{figure}[tbp]
\centering
\begin{tikzpicture}[x=1cm,y=1cm,font=\footnotesize,
  forcing label/.style={align=left,inner sep=2pt},
  forcing tube/.style={line width=14pt,line join=round,line cap=round},
  forcing wall/.style={draw=orange!45,line width=20pt,
    line join=round,line cap=round},
  forcing core/.style={draw=teal!70!black,line width=4pt,
    line join=round,line cap=butt}]
  \coordinate (Xzero) at (140:3.8);
  \coordinate (Xone) at (302:3.8);
  \coordinate (Rzero) at (-2.60,2.17);
  \coordinate (Rone) at (1.791,-2.867);
  \coordinate (Szero) at (-.38,1.20);
  \coordinate (Sone) at (1.97,1.20);

  \draw[draw=orange!45,line width=5pt] (0,0) circle (3.55);
  \draw[forcing wall] (Rzero) -- (-1.3,2.17) -- (-.85,1.70)
    -- (-.75,1.20) -- (Szero);
  \draw[forcing wall] (Rone) -- (1.791,-1.7) -- (2.65,-1.05)
    -- (2.65,1.20) -- (Sone);
  \filldraw[fill=white,draw=orange!45,line width=5pt]
    (-.50,.73) rectangle (2.10,1.70);
  \foreach \p in {Xzero,Xone} {\fill[white] (\p) circle (.49);}
  \draw[forcing tube,white] (Xzero) -- (Rzero) -- (-1.3,2.17)
    -- (-.85,1.70) -- (-.75,1.20) -- (Szero);
  \draw[forcing tube,white] (Xone) -- (Rone) -- (1.791,-1.7)
    -- (2.65,-1.05) -- (2.65,1.20) -- (Sone);
  \fill[white] (-.46,.77) rectangle (2.06,1.66);
  \draw[forcing tube,draw=blue!9] (Rzero) -- (-1.3,2.17)
    -- (-.85,1.70) -- (-.75,1.20) -- (Szero);
  \draw[forcing tube,draw=blue!9] (Rone) -- (1.791,-1.7)
    -- (2.65,-1.05) -- (2.65,1.20) -- (Sone);

  \draw[dashed,blue!65!black,line width=.7pt]
    (-.55,.60) rectangle (2.15,1.83);
  \node[anchor=south east,text=blue!55!black] at (2.15,1.87) {$\Omega$};
  \filldraw[fill=green!10,draw=green!45!black,line width=.7pt]
    (-.20,.85) rectangle (1.80,1.55);
  \node[forcing label,align=center] at (.80,1.28) {$V$};
  \foreach \lo/\hi/\slo/\shi in {-.48/.14/-.42/.05,1.56/2.08/1.63/2.04} {
    \draw[blue!55!black,densely dashed] (\lo,.94) rectangle (\hi,1.46);
    \draw[blue!45!black] (\slo,1.05) rectangle (\shi,1.35);
  }

  \fill[orange!25] (-1.98,1.925) rectangle (-1.50,2.415);
  \fill[orange!25] (2.405,-.84) rectangle (2.895,-.35);
  \draw[forcing core] (Rzero) -- (-1.3,2.17) -- (-.85,1.70)
    -- (-.75,1.20) -- (Szero);
  \draw[forcing core] (Rone) -- (1.791,-1.7) -- (2.65,-1.05)
    -- (2.65,1.20) -- (Sone);
  \draw[black!60,densely dashed] (-1.98,1.925) -- (-1.98,2.415)
    (-1.50,1.925) -- (-1.50,2.415)
    (2.405,-.84) -- (2.895,-.84)
    (2.405,-.35) -- (2.895,-.35);
  \node[anchor=south] at (-1.74,2.48) {marker};
  \node[anchor=east] at (2.33,-.595) {marker};
  \node[anchor=east,text=blue!60!black] at (-1.08,1.23) {$T_0$};
  \node[anchor=east,text=blue!60!black] at (1.51,-2.06) {$T_1$};

  \draw[black!60,densely dotted] (Xzero) -- (Rzero)
    (Xone) -- (Rone) (Szero) -- (-.10,1.20) (Sone) -- (1.70,1.20);
  \draw[black!75,dashed] (-.10,1.20) -- (.20,1.02)
    -- (1.30,1.02) -- (1.70,1.20);
  \foreach \p in {Rzero,Rone,Szero,Sone} {
    \fill[teal!70!black] (\p) circle (1.2pt);
  }
  \fill[black] (-.10,1.20) circle (1.1pt);
  \fill[black] (1.70,1.20) circle (1.1pt);
  \node[anchor=south west,inner sep=1pt] at (.02,1.22) {$q_0$};
  \node[anchor=south east,inner sep=1pt] at (1.58,1.22) {$q_1$};
  \node[anchor=north east,text=blue!55!black] at (-.83,.52) {$G_0^{\pm}$};
  \draw[blue!45!black] (-.79,.52) -- (-.34,.94);
  \node[anchor=north,text=blue!55!black] at (1.95,.43) {$G_1^{\pm}$};
  \draw[blue!45!black] (1.95,.46) -- (1.95,.94);

  \draw[black!75,line width=.8pt] (0,0) circle (3.8);
  \node[anchor=south] at (0,3.84) {$\partial K$};
  \foreach \p in {Xzero,Xone} {
    \draw[blue!55!black,densely dashed] (\p) circle (.62);
    \draw[blue!45!black] (\p) circle (.46);
    \fill[black] (\p) circle (1.6pt);
  }
  \node[anchor=east] at ($(Xzero)+(-.76,.06)$) {$x_0$};
  \draw[black!50] ($(Xzero)+(-.12,.01)$) -- ($(Xzero)+(-.70,.06)$);
  \node[anchor=west] at ($(Xone)+(.76,-.12)$) {$x_1$};
  \draw[black!50] ($(Xone)+(.12,-.02)$) -- ($(Xone)+(.70,-.12)$);
  \node[anchor=south,text=blue!55!black] at ($(Xzero)+(.24,.65)$) {$A_0^{\pm}$};
  \node[anchor=north,text=blue!55!black] at ($(Xone)+(-.25,-.64)$) {$A_1^{\pm}$};
  \filldraw[fill=white,draw=black!65] (0,0) circle (.2375);
  \node at (0,0) {$I$};
  \node[forcing label,anchor=east] at (-.70,-.45) {untouched};
  \draw[black!50] (-.66,-.43) -- (-.19,-.13);

  \filldraw[fill=green!10,draw=green!45!black]
    (4.25,2.30) rectangle (4.90,2.54);
  \node[forcing label,anchor=west] at (5.04,2.42) {Preserved patch\\$f=b$};
  \draw[forcing core] (4.25,1.40) -- (4.90,1.40);
  \node[forcing label,anchor=west] at (5.04,1.40) {Negative cores\\$f=-J$};
  \draw[draw=orange!45,line width=5pt] (4.25,.38) -- (4.90,.38);
  \node[forcing label,anchor=west] at (5.04,.38) {Wall band $W_*$\\$f=H$};
  \draw[dashed,blue!65!black] (4.25,-.70) -- (4.90,-.70);
  \node[forcing label,anchor=west] at (5.04,-.70) {Denominator\\boundary};
  \node[forcing label,anchor=north west,font=\scriptsize,text=black!65]
    at (4.25,-1.47) {Schematic; smooth\\transition collars omitted.};
\end{tikzpicture}
\caption{The fixed geometry in Lemma~\ref{comp:finite-catalogue-construction}.
The preservation patch $V$ and its terminal caps lie inside the denominator
$\Omega$, away from the untouched disk $I$.  Two disjoint polygonal tubes
connect the outer caps at $x_0,x_1\in\partial K$ to opposite faces of $V$.
Their negative cores stop inside the safe caps; the dashed preserved
certificate joins $q_0$ to $q_1$ and remains wholly in $V$.  For $D_f=e^{\xi f}\cdot D$, the wall
band prevents a path of $D_f$-length at most $B$ from crossing the forbidden
gap.  The two marked
cross-sections force every such complete tube passage to approach its
core: the marker has height $H$ outside the core collar.  Exact offset
boundaries and smoothing collars are suppressed.}
\label{fig:forcing-corridors}
\end{figure}


\begin{lemma}[Construction of the finite forcing catalogue]
\label{comp:finite-catalogue-construction}
Let $K=\overline{B_2(0)}$ and
$I=\overline{B_{1/8}(0)}$.  Fix a finite set of Euclidean
configurations of the following type.
\begin{enumerate}[label=(\alph*)]
\item A rectangular preservation patch $V$ and a larger
      rectangular denominator domain $\Omega$ satisfy
      $\overline V\subset\Omega\Subset
      B_{3/2}(0)\setminus\overline{B_{1/4}(0)}$.
      Two endpoint ports lie near opposite faces of the
      rectangle inside~$V$.
\item Disjoint terminal caps
      $G_j^-\Subset G_j^+\Subset\Omega$, $j=0,1$, contain
      those ports.  Each $G_j^-$ also contains a fixed stop
      point $s_j\notin\overline V$, which can approach the
      corresponding face from its outer side.  The caps have
      been chosen by a sufficiently fine finite gridding of
      the rectangle and its two endpoint ports.
\item Disjoint outer caps $A_j^-\Subset A_j^+$ surround two
      separated gridded points of $\partial K$.  A fixed
      point $r_j\in A_j^-\cap\operatorname{int}K$ lies a
      positive distance inside $K$.  The possible exterior
      attachment points $x_j\in\partial K$ lie in a compact
      subset of $A_j^-$.
\end{enumerate}
All indicated containments and separations have strict
Euclidean margins.  The endpoint ports and the terminal stops
can, for example, be obtained by taking a thin collar of a
gridded rectangle and cap boxes straddling its opposite faces.
Fix a finite set $\mathcal B\subset\mathbb R$ of exact
constant shifts.

For each configuration one can choose finitely many fixed
polygonal core domains and marker rectangles as described
below.  Suppose the reference metric $D$ and the chosen
endpoints $q_j$ satisfy the following deterministic bounds,
uniformly for that finite family, with positive numbers
$u,w,e,A,K_*,m_*$ and a modulus $\omega(t)\to0$.
\begin{enumerate}[label=(\roman*)]
\item Every point of $A_j^+$ has an old local connection to
      $x_j$ of length at most $e$; there is a connection from
      $x_j$ to $r_j$ in $A_j^-$ of length at most $e/3$.
\item Every point of $G_j^+$ has a $D_b$-connection to $q_j$
      in $\Omega$ of length at most $e$, where
      $D_b=e^{\xi b}D$.  There is a connection from $s_j$ to
      $q_j$ in $G_j^-$ of $D_b$-length at most $e/3$.
\item There is a $D_b$-connection from $q_0$ to $q_1$ inside
      $V$ of length at most $Au$.
\item Each point of the closed core domain can be connected
      to its terminal stop, within that core, at old cost at
      most $K_*w$.  The outer stop has such a connection too.
\item The fixed forbidden-gap and marker crossings have old
      length at least $m_*w$.  Points at Euclidean distance
      at most $t$ from a core can be connected to that core
      inside its prescribed open collar at old cost at most
      $\omega(t)w$.
\end{enumerate}
Fix a local budget $B$ with $Au+2e\le B$.  Then there is a
finite deterministic list of smooth functions $f$, depending
only on the geometric list, these bounds, and $\mathcal B$,
with the following properties.
Each function is supported compactly in
$B_2(0)\setminus I$, is exactly $b$ on $V$, and satisfies
the network, gate, and cheap-approach hypotheses of
Lemma~\ref{comp:gate-forcing} for the indicated configuration
and budget $B$, with all marker-to-terminal joins and the preserved
certificate compactly confined to the bounded neighborhood
$W=B_3(0)$.  It also provides a competitor of cost at
most $\ell_0+\ell_1+B$ whenever the exterior approaches to
$x_0,x_1$ have old costs $\ell_0,\ell_1$.
The list has uniformly bounded amplitudes and Dirichlet
energies, including after any Euclidean translation and
dilation.
\end{lemma}

\begin{proof}
We first make all Euclidean choices; no metric-dependent
curve is used to choose the shape of a shift.

\paragraph{Two embedded corridors.}
Extend the short prescribed initial and terminal pieces to
an arc joining $\partial K$ to the corresponding face of
$\partial V$, inside
$\operatorname{int}K\setminus(I\cup\overline V)$.
Choose the first arc simple.  An arc joining distinct
boundary components of this planar region does not
disconnect its interior after cutting along the arc: the
outer boundary and the rectangular boundary become one
boundary component, while the omitted disk remains another.
Thus the second pair of prescribed ends can be joined in
the cut region.  Approximate both arcs by simple polygonal
arcs, keeping the prescribed short end pieces and positive
clearance from each other, the omitted disk, and the
rectangle except in their terminal caps.  Trim their ends
to obtain disjoint arcs from $r_j$ to $s_j$.

One can require the arcs to avoid the other attachment and
terminal cap by choosing those caps small first.  At a
terminal cap approach along the outward normal to its
specified face.  Since the input configurations form a
finite list, these choices can be made once for each member.
Their minimum positive clearance is then positive.

\paragraph{Nested neighborhoods and the region graph.}
Choose three nested regular neighborhoods
\[
 O^0\subset O^1\subset O^2,\qquad
 N^0\Subset N^1\Subset N^2,\qquad
 T_j^0\Subset T_j^1\Subset T_j^2,
\]
where $O^k$ contains the exterior of $K$ and a thin inner
collar, $N^k$ surrounds $\overline V$, and $T_j^k$ is a
regular tube around the $j$th polygonal arc with end bulbs.
For the unbounded sets $O^k$, require positive separation
between their finite boundary components.
Choose the neighborhoods sufficiently thin that
\[
 \overline{N^2}\subset\Omega,\quad
 O^2\cap N^2=\varnothing,\quad
 T_0^2\cap T_1^2=\varnothing,
\]
and their only overlaps satisfy
\[
 O^2\cap T_j^2\subset A_j^+,
 \qquad N^2\cap T_j^2\subset G_j^+.
\]
All these incidences persist under a further small thickening.
This follows directly from the positive clearances of the
arcs away from their four prescribed end caps.  Make the
end bulbs large enough that the smaller safe caps
$\overline{A_j^-}$ and $\overline{G_j^-}$ lie in the union
of the smallest neighborhoods, with positive margin.
There is no requirement that the larger effective caps be
cheap in the modified metric.

Write
\[
 S^k=O^k\cup N^k\cup T_0^k\cup T_1^k.
\]
These sets avoid a fixed neighborhood of $I$, and there is
$\rho_*>0$
such that the Euclidean $\rho_*$-neighborhood of
$\overline{S^0}$ lies in $S^2$.
The middle layer
\[
 W_*=
 \{z:\rho_*/3\le\operatorname{dist}(z,S^0)
                         \le2\rho_*/3\}
\]
is compactly contained in
$\operatorname{int}K\setminus I$ and disjoint from all
the safe caps and the preservation patch.  A curve starting
in $S^0$ and leaving $S^2$ has a subcurve in $W_*$ whose
endpoints are separated by at least $\rho_*/3$.
This is the forbidden-gap crossing in (v).

The nerve of the cover of $S^2$ by its four regions is
precisely the cycle
\[
 O^2-T_0^2-N^2-T_1^2-O^2.
\]
Any compact curve contained in $S^2$ admits a finite
subdivision subordinate to this open cover and hence a
finite itinerary in that graph.  This statement allows
arbitrarily many incidental boundary contacts.

\paragraph{Core domains and marked sections.}
Choose closed polygonal core domains $C_j\subset T_j^0$
around the trimmed arcs, with the following margins:
\[
 C_j\cap O^2=\varnothing,\qquad
 C_j\cap\overline V=\varnothing,
\]
and both stops $r_j,s_j$ lie in $C_j$.
The outer collar is chosen thinner than the depth of $r_j$;
the inner stop lies outside $\overline V$ by assumption.
The core connection bounds in (iv) concern these fixed
closed domains.  They may be obtained from the internal
box estimates on a finite rectangular cover, including
access to the flat faces.

On a straight interior piece of each polygonal tube, away
from its end bulbs and from $O^2\cup N^2$, select a closed
longitudinal marker rectangle spanning the full closed transverse
section of $T_j^2$ between its two end cross-sections.
The tube there is an
ordinary Euclidean strip.  A path in $T_j^2$ from its outer
overlap to its terminal overlap must cross this rectangle
between two cross-sections a fixed positive distance apart.
One can see this by using the longitudinal coordinate in
the regular tube; retain the subcurve between its last
visit to the first cross-section and its first subsequent
visit to the second.  This is the marker crossing in (v).

\paragraph{A smooth wall and constant plateaux.}
Choose $H>\max\{0,\max\mathcal B\}+1$ so large that
\[
 e^{\xi H}m_*w>B.
\]
Choose a smooth baseline $a$, compactly supported in
$\operatorname{int}K\setminus I$, which is zero on the
safe outer caps and near $\partial K$, equals $b$ on a
neighborhood of $\overline V\cup\overline{G_0^-}
\cup\overline{G_1^-}$, and equals $H$ on a neighborhood
of $W_*$ and on both marker rectangles.  Require also
\[
 a\ge0\text{ on }O^2,
 \qquad a\ge b\text{ on }N^2,
 \qquad \min\{0,b\}\le a\le H.
\]
These requirements have disjoint closed cores and positive
margins.  Explicitly, first choose a cutoff $\zeta$ equal
to one on the preservation patch and the safe terminal
caps, supported in $S^0$ and away from $O^2$, the marker rectangles,
the wall layer, and the safe outer caps.  Next choose
$\psi$ equal to one on $N^2$, the support of $\zeta$,
the marker rectangles, and $W_*$, and zero on the safe
outer caps and near the forbidden support boundaries.
The formula
\[
 a=H\psi+(b-H)\zeta
\]
has all the stated properties, with the cutoffs chosen
between zero and one.

Choose $J>\max\{0,-\min\mathcal B\}+1$ so that
\[
 e^{-\xi J}K_*w<e/3.
\]
Then choose $\rho>0$ smaller than all the fixed geometric
margins, and so small that
\[
 e^{\xi H}\omega(\rho)w<e/3.
\]
Let $\chi_j$ be smooth, equal to one on $C_j$, and supported
in its $\rho$-neighborhood.  These supports are disjoint,
lie inside $S^0$, and avoid $O^2\cup\overline V\cup W_*$.
Set $\chi=\chi_0+\chi_1$ and define
\[
 f=(1-\chi)a-J\chi.
\]
This is a smooth compactly supported shift.  It equals
$-J$ on the cores, $b$ on $V$, and $H$ on $W_*$.  On
the safe terminal caps $f\le b$, and on the safe outer
caps $f\le0$.  Moreover,
\[
 f\ge0\text{ on }O^2,
 \qquad
 f\ge b\text{ on }N^2\setminus(T_0^2\cup T_1^2).
\]

\paragraph{The local budget gives the itinerary.}
A curve starting at $\partial K$ starts in $S^0$.
If a curve of total new length at most $B$ leaves $S^2$,
it crosses $W_*$ and pays more than $B$ by the choice
of $H$.  Thus every local-budget portion between the first
and last visits to $\partial K$ stays in $S^2$.  The
finite-cover argument above gives its claimed itinerary.
The omitted disk supplies no bypass, since it lies outside
$S^2$.  The outer and terminal gate-error conditions follow
from their containments in $A_j^+$ and $G_j^+$ and (i)--(ii).

\paragraph{A marked passage gives cheap endpoint access.}
On a marker rectangle, $a=H$.  Outside the
$\rho$-neighborhood of its core one has $\chi=0$ and
therefore $f=H$.  A complete corridor passage avoiding
that core neighborhood would have an all-wall marker
crossing and cost more than $B$.  Every complete passage
with the local budget consequently contains a point $z$
at distance at most $\rho$ from the core.  By (v), this
point connects to the core at new cost at most
$e^{\xi H}\omega(\rho)w<e/3$.  By (iv), it then reaches
$s_j$ through the core at cost less than $e/3$.  The safe
terminal cap has shift at most $b$, so (ii) gives a final
connection to $q_j$ at cost at most $e/3$.  Thus
\[
 D_f(z,q_j;\text{prescribed neighborhood})\le e.
\]
These marker points supply cheap approach joins.  The
different effective gate contacts are used only for the
old $D_b$ lower-bound corrections in the forcing lemma.
The prescribed core collars and terminal-cap paths lie inside
$B_2(0)$, and the retained certificate lies inside $V$.
Consequently the shortcut assembled from these joins and that
certificate has compact confinement in $W=B_3(0)$.  The outer
region and the long anchor approaches play no part in this
confinement assertion.

\paragraph{The competitor and finiteness.}
Follow the unchanged exterior approaches, the safe outer
cap connections, the negative cores, the safe terminal
cap connections, and the connection through $V$ from (iii).
The six approach contributions have total cost at most
$2e$, so the resulting competitor has length at most
$\ell_0+\ell_1+Au+2e\le\ell_0+\ell_1+B$.

Every region, core, marker, cutoff, and height was selected
from deterministic data after the finite geometric list
and uniform metric bounds had been fixed.  Hence these
functions form a finite list.  Their amplitudes are
bounded by $\max\{H,J,\max_{b\in\mathcal B}|b|\}$.
Smoothness and compact support give a finite common
Dirichlet-energy bound.  Finally,
\[
 \int_{\mathbb C}
 \left|\nabla f\left(\frac{z-z_0}{r}\right)\right|^2
 \frac{dz}{r^2}
 =\int_{\mathbb C}|\nabla f(w)|^2\,dw,
\]
so the energy of $z\mapsto f((z-z_0)/r)$ is unchanged
by translation and positive dilation.
\end{proof}

\begin{remark}[Which bounds are chosen when]
\label{comp:construction-choice-order}
The preservation patch, terminal caps, and their error
bounds are chosen first, followed by the thin-domain
connection bound $A$.  Next choose the fixed ratio $w/u$
large enough for the exterior gap and budget conditions.
The outer attachment grid and all Euclidean tubes are then
fixed.  The new core and gap bounds $K_*,m_*,\omega$ are
truncated on those finitely many domains.  These later
truncations only affect $H,J,\rho$; they do not alter $A$
or the terminal caps.  Thus the construction has no
circular dependence of the upper compartment budget on
the later routing constants.
\end{remark}

\begin{lemma}[Uniform stability for a shrinking modification]
\label{comp:shrinking-modification-stability}
Let $D$ be a proper geodesic metric inducing the Euclidean
topology and having a unique geodesic $P$ between fixed distinct
anchors $a_0,a_1$, parametrized by fractions of its length.
Fix $0<t_1<t_2<1$ and an open set $U$ containing
$P([t_1,t_2])$ with positive clearance.  Suppose the old passage
rule has upper comparison $C$ in open subsets of $U$ and
\[
 F_U(P(t_1),P(t_2))
 =C(t_2-t_1)D(a_0,a_1).
\]
Consider bounded shifts $f$ with one common amplitude bound,
supported in balls whose radii tend uniformly to zero and whose
centers lie at vanishing distance from a fixed compact middle
subarc of $P((t_1,t_2))$.
Let $F'$ be any output passage rule with upper comparison $C$
against $D_f=e^{\xi f}\cdot D$.  Suppose it preserves all old
exterior certificate restrictions outside a slightly enlarged
support ball.  Assume the certificate system permits marked
subpath extraction with additive nonnegative costs and
concatenation with arbitrarily small open loss.

Uniformly over these modifications and all output fillings,
every fractionally parametrized $D_f$-geodesic $P'$ satisfies
\[
 \sup_{t\in[0,1]}D(P'(t),P(t))\longrightarrow0,
 \qquad
 D_f(a_0,a_1)\longrightarrow D(a_0,a_1),
\]
and
\[
 \frac{F'_U(P'(t_1),P'(t_2))}
 {(t_2-t_1)D_f(a_0,a_1)}\ge C-o(1).
\]
If a strict shortcut along $P'$ is contained in the shrinking
modification neighborhood, then for all sufficiently small
supports this ratio is strictly less than $C$.
\end{lemma}

\begin{proof}
Choose constants $0<m\le1\le M<\infty$ with
$mD\le D_f\le MD$.  Properness puts all the relevant geodesics
in one compact $D$-ball.  Let $\omega_n$ be the supremum of the
old $D$-diameters of suitably enlarged possible support balls.
The Euclidean radii tend uniformly to zero, so continuity gives
$\omega_n\to0$.

Replace the portion of an old geodesic between its first and
last support visits by an old minimizing connection.  Its new
length is at most $M\omega_n$.  For the reverse comparison,
cut a new geodesic at its first and last support visits and
insert an old minimizing connection of cost at most
$\omega_n$.  Outside the support the lengths coincide.  Thus
\[
 |D_f(x,y)-D(x,y)|\le(M+1)\omega_n
\]
for the endpoints under consideration.  Moreover, the new
geodesic segment spanning all its support visits has new length
at most $M\omega_n$, and old length at most
$m^{-1}M\omega_n$.  Hence the old length of the entire new
geodesic tends to $D(a_0,a_1)$.

The new geodesics are uniformly Lipschitz in $D$ under their
fractional parametrizations.  Every subsequential limit is an
old minimizing path; equivalently, the uniform distance
comparison shows that each of its subsegments has the
appropriate fractional distance.  Uniqueness identifies the
limit with $P$ with its specified parametrization.  A
contradiction-subsequence argument proves uniformity over all
support locations and bounded shifts.  Write the uniform
endpoint discrepancy as $\eta_n\to0$.

Put $q_0=P(t_1)$, $q_1=P(t_2)$ and
$q'_0=P'(t_1)$, $q'_1=P'(t_2)$.  Consider an arbitrary
output certificate from $q'_0$ to $q'_1$ compactly in $U$.
If it enters the modification neighborhood, cut its
approximants at their first and last visits to a ball strictly
larger than the ball removed from the preserved exterior
patch.  Marked subpath extraction gives retained prefix and
suffix certificates in that open exterior patch, with total
cost no larger than the original cost.  Transfer them to the
old sample.  Join their cut points at old passage cost at
most $C\omega_n+o(1)$, and correct the two endpoints at total
cost at most $2C\eta_n+o(1)$.  All corrections lie in $U$,
by its clearance and the fact that short old geodesics remain
in small neighborhoods.  For a certificate avoiding the
larger ball, transfer it directly and make only the endpoint
corrections.  The retained pieces use their own positive
clearance from $\partial U$; a common clearance for all
competing certificates is unnecessary.

Taking infima yields, with
$\ell=(t_2-t_1)D(a_0,a_1)$,
\[
 C\ell\le F'_U(q'_0,q'_1)+C(\omega_n+2\eta_n).
\]
The new denominator $\ell'$ tends to $\ell>0$, proving the
ratio bound.  Finally, the old middle subarc is separated
from $P([0,t_1]\cup[t_2,1])$.  Uniform convergence therefore
places every visit of $P'$ to the shrinking modification
neighborhood strictly between $t_1$ and $t_2$.  Concatenating
the strict shortcut with upper-comparison paths along the
remaining geodesic portions gives the strict upper inequality.
\end{proof}

\begin{remark}[Deterministic errors before counting]
\label{comp:deterministic-stability-errors}
The error in Lemma~\ref{comp:shrinking-modification-stability}
is a random function of the incoming sample, uniform over all
candidate modifications.  Its almost sure convergence to zero
allows a deterministic sequence $\epsilon_N\to0$ for which
the error exceeds $\epsilon_N$ with probability tending to
zero.  Intersect the event that the error is at most $\epsilon_N$ with
the positive-probability saturation event and the every-point
good-cover event before
forming the source submeasures.  A positive lower cutoff for
the interval length ensures that the ratio error is controlled.
This gives the deterministic common-law sets $H_N$ required
by Lemma~\ref{comp:weighted-output-multiplicity}.
\end{remark}

It remains to make a successful catalogue available with high
probability in a reference shell.  When the passage kernel uses the
genuine order of the local imaginary traversals, an independent trial
must determine that order from its own dependency disk.  The next lemma
provides such disks; for an order-free kernel this step is unnecessary.

\begin{lemma}[A locally certified imaginary order]
\label{comp:local-order-certificate}
Let $\widehat h$ be the whole-plane imaginary field modulo its
angular period, and let $\eta$ be its Lebesgue-parameterized
space-filling curve.  Fix a compact set $K\Subset U$, with $U$
bounded and open.  The following public properties suffice:
\begin{enumerate}[label=(\roman*)]
\item the rational-point visit times of $\eta$ are dense;
\item its full directed $U$-traversals indexed by rational points,
      parametrized by Lebesgue area from their entry times,
      are determined from $\widehat h$ on an arbitrarily small
      enlargement of $U$, and the same holds for bounded disks;
\item each rational spatial point is visited once almost surely.
\end{enumerate}
There are events $D_R$, measurable from the imaginary field on
a slightly enlarged ball of radius $R$, such that
$\mathbb P(D_R)\to1$ and, on $D_R$, the genuine order of all
$U$-traversals meeting $K$ is determined by that restriction.
\end{lemma}

\begin{proof}
Transience confines the visits to $\overline U$ to a compact
time interval.  Uniform continuity gives a positive lower
duration for each $U$-traversal reaching $K$, because it must
cross $\operatorname{dist}(K,\partial U)>0$.  There are
therefore only finitely many such traversals.  Property (i)
ensures that each contains a rational-point visit.  In the
locally reconstructed collection select the first rational
representative in a fixed enumeration for each traversal;
this selection does not use their chronological order.

Apply (ii) again to the disk $B_R$.  For every selected rational
representative, reconstruct its full directed $B_R$-traversal,
with time zero at entry.  Let $D_R$ be the event that all these
records are identical.  It is measurable from the buffered
imaginary-field restriction.  Distinct traversal intervals cannot
have identical records: they have positive time length and hence
positive area, whereas ranges of disjoint time intervals overlap
in zero area.  On $D_R$, the common directed parametrization and
(iii) specify the order of the representatives exactly, and therefore
the order of their distinct small-domain traversals.

Almost surely the finitely many representative visit times have a
finite minimum and maximum.  By continuity, the image of the
interval between them is compact and lies in $B_R$ for all
sufficiently large $R$.  They then lie in one common $B_R$-traversal,
so $D_R$ holds eventually.  This proves the exhaustion.
\end{proof}

\begin{remark}[Public inputs for the preceding lemma]
Properties (i) and (ii) are given in Section~2.1.3 and
Lemma~2.4 of Gwynne--Miller--Sheffield
\cite{GwynneMillerSheffieldHarmonic}.
For (iii), area parametrization makes the set of multiply visited
spatial points Lebesgue-null: cover pairs of distinct visit times
by disjoint rational time intervals, whose ranges have zero-area
intersection.  Translation invariance of the whole-plane curve
modulo time translation then gives zero probability that any fixed
spatial point is multiply visited, and a countable intersection
handles rational points.  The stationary angular-phase convention
for the imaginary field is used here.
The argument uses density of rational \emph{visit times};
positive area of an excursion range alone would not imply
that it contains a rational point.
\end{remark}

\begin{proposition}[Conditional availability of a finite catalogue]
\label{comp:conditional-catalogue-availability}
Consider favorable normalization indices $v$ and a family of
passage kernels.  Suppose:
\begin{enumerate}[label=(\roman*)]
\item For each open preservation patch $V$, the restricted
      certificate kernel is a function of $h|_V$, the local
      imaginary traversal data, and their finite genuine
      order.  The order depends only on the imaginary field,
      and the kernels are compatible under restriction to
      smaller open patches.  Separated restrictions have
      conditional product kernels.
\item Under a common normalized background-field law, at every
      favorable index there is probability at least $\beta>0$
      of a strict low-ratio certificate whose entire range lies in
      a retained closed rectangular envelope and whose endpoints
      lie on its opposite faces, with
      common geometric slack and a denominator in a fixed
      buffered domain.  The fixed reference-metric moduli and
      finite-domain connection bounds are almost surely finite.
\item The exact affine and volume-normalization identities of
      the joint kernel family transport each trial to the
      corresponding favorable reference index after its real
      coarse constant, spatial logarithm, and fixed volume
      displacement have been compensated.  These are identities
      of transported fields and deterministically normalized
      observations, not a Weyl rule for a fixed passage object.
\item The normalized GFF harmonic profiles, after removing the
      real coarse constant and reducing the imaginary constant
      modulo its period, have uniformly tight interior smooth
      norms on fixed-ratio buffered trial disks.
\end{enumerate}
For every $\theta>0$, one can choose finitely many disjoint
buffered trials, a finite deterministic grid of exact real
shifts, and the metric bounds for a finite forcing catalogue
so that, uniformly over favorable indices, the completed
catalogue is available with probability at least $1-\theta$.
The good event can include fixed bounds on all the resulting
Cameron--Martin pairings.  No continuity of passage costs under
rounding a field shift is required.
\end{proposition}

\begin{proof}
First truncate the normalized small-box and face modulus at
a loss smaller than $\beta/10$.  Choose a small mesh width $t$
using this common internal box modulus and the prescribed terminal
error budget.  Grid both the closed-envelope parameters and the
positions of the two endpoints.  For each resulting finite bin,
choose an outward rectangular preservation collar $V$ containing
every envelope in the bin with positive margin.  Each indicated
face of $V$ is within $O(t)$ of the corresponding envelope face.
Choose terminal caps of diameter $O(t)$ straddling these two faces,
containing the endpoint bins and fixed stop points outside
$\overline V$.  The internal box modulus on their collars gives
the desired arbitrarily small terminal connection errors.
The entire certificate range is in $V$ because it was retained
inside its closed envelope.  Choose an intermediate denominator
domain containing $\overline V$ inside the original one.
Its intrinsic denominator is at least the old denominator, so the
strict low-ratio bound remains valid.

The possible endpoints lie in deterministic compact subsets
of finitely many newly thin domains.  Their internal reference
connection diameters are finite almost surely.  Fix an upper
bound $A'$ with a further loss smaller than $\beta/10$.
These cap sizes, thin domains, and $A'$ remain fixed during
all later tube truncations.

For a singleton order input, take its determination event to be
the whole sample space.  Otherwise choose a compact $K\Subset U$
covering the prepared local traversals.  By
Lemma~\ref{comp:local-order-certificate}, a fixed imaginary buffer
makes the sum of order-determination failure probabilities over
the finite geometry bins smaller than $\beta/10$.
Write $E_{v,a}$ for the prepared success event in
bin $a$ at favorable index $v$, intersected with its determination
event, and put $E_v=\bigcup_a E_{v,a}$.  The genuine order on each
of these events is a function of the buffered imaginary field.
Thus the success union $E_v$ is local and retains a common positive
probability, say $\beta_1$.  No individual bin is required to have
positive probability.

Choose a larger fixed conditioning disk around the complete
dependency region.  Local absolute continuity transfers the
reference fields to a pair of Dirichlet GFFs on this disk,
restricted to its interior.  The field density is independent
of the index.  Truncate this single density to conclude that
the success union $E_v$ has probability at least
$\beta_2>0$ under this fixed Dirichlet reference law, uniformly
in the index.

For an affine trial copy, condition on the fields outside its
conditioning disk.  The fields inside are Dirichlet fields
plus exterior-measurable harmonic functions.  On an acceptable
profile event, subtract the real coarse constant and reduce
the imaginary constant modulo its period.  Smooth cutoffs
extend the remaining profiles on the dependency region to
Cameron--Martin functions of uniformly bounded energy.
The same remains true after adding any real residual constant
in~$[-1,1]$.

Let $P$ be the fixed Dirichlet reference law, including an
independent auxiliary kernel seed, and $Q$ one such
conditional law.  The passage kernel is the same function of
the field inputs under both laws.  A uniform inverse
likelihood second moment gives, for the success union $E_v$,
\[
 \beta_2^2
 \le Q(E_v)\,\mathbb E_Q\left[
           \left(\frac{dP}{dQ}\right)^2\right],
 \qquad Q(E_v)\ge\beta_*:=\beta_2^2/K_*>0.
\]
The profile cutoff and $\beta_*$ are independent of the number
of trial copies: the normalized domains and buffer ratios are
fixed, and (iv) controls their harmonic oscillations uniformly.

Choose that cutoff so the expected fraction of unacceptable
profiles is below $\theta/100$.  Markov's inequality implies
that at least half the trials are acceptable except with
probability at most $\theta/50$.  Choose $J$ with
\[
 (1-\beta_*)^{J/2}<\theta/10,
\]
and pack $J$ deterministic disjoint copies of the entire
conditioning disk into the prescribed shell subregion.
All spatial ratios are now fixed.  Before choosing $M$, allocate a
fixed portion of the remaining failure probability to the outer
reference shell: fix a positive normalized crossing lower bound
$a$, a finite boundary-connection upper bound, and a reference
connection modulus there.  The normalized reference field law and
shell geometry are fixed, so these choices do not depend on $M$.
Write $\varepsilon$ for the already fixed normalized terminal error
budget.  Choose the fixed quantum separation $M$ required by the forcing
budget, using this $a$, the already fixed $A'$ and cap errors;
in particular impose $A'+4\varepsilon<aM^p$ in the units of
Lemma~\ref{common:forcing-verification}.  Assumption (iii) absorbs
its logarithmic displacement into the desired compensation,
so it does not change $\beta_*$.

The $J$ desired compensations are exterior-measurable.  Their
normalized vector is uniformly tight: the favorable index
cancels in (iii), leaving only the fixed affine charts, their
coarse field values, and the now fixed displacement from $M$.
Truncate this vector to a compact interval with loss at most
$\theta/10$, and cover that interval by a finite grid
$\mathcal B$ of mesh at most two.  Select the nearest grid
point using only exterior data.  Its residual lies in
$[-1,1]$, precisely the range in the preceding likelihood
bound.

For every exact $b\in\mathcal B$, draw a fresh auxiliary
observation array for the prepared patches from the common kernel
at the exactly shifted real field, and restrict it to each $V$.
Use independent
seeds across disks and grid values, conditional on the fields.
For the selected grid point declare success on the union of its
events $E_{v,a}$; select the first successful bin in the fixed
finite ordering.  On that bin's determination event the genuine
order equals the locally computed order.  Hence the success indicator is
a function of that disk's field restriction, its independent
seed, and the exterior-measurable grid choice.  The indicators
are conditionally independent given the exterior of the
disjoint disks.  Each acceptable disk has success probability
at least $\beta_*$, and the chance of no success is at most
$(1-\beta_*)^{J/2}$.

The successful certificate is kept at the exact grid value
at which it was drawn.  Rounding has been handled by changing
the field law in a uniform absolute-continuity estimate;
no claim about persistence of an individual certificate under
rounding has been made.

Finally impose the remaining regularity and finite internal bounds
for all the gridded routing, core, gap, and access domains, with total additional
failure smaller than the unused portion of $\theta$.
These later bounds affect only the wall and core heights and
the smoothing widths.  They do not change the cap modulus
or $A'$.  After the resulting finite smooth recipe list is
fixed, truncate its finitely many Gaussian linear pairings.
For a smooth compactly supported recipe, this linear functional is a
distributional pairing against its Laplacian, so it depends only on
the real field on the recipe support.  Clipping is therefore part of
the same local shell event.  The uniformly bounded energies give a
fixed clipping threshold with arbitrarily small failure.  Allocating the
remaining losses proves the claimed probability.
\end{proof}

\begin{remark}[The real and imaginary dependency patches differ]
\label{comp:separate-field-dependencies}
The larger order buffer in the preceding proof is an
\emph{imaginary}-field dependency region.  A real-field
modification leaves this input unchanged.  By the restriction
compatibility in assumption (i), a retained certificate
compactly in $V$ uses the real field only on $V$, even if it
was originally sampled as part of a larger auxiliary passage
object.  Thus a recipe equal to $b$ on $V$ retains precisely
the required kernel after modifying the real field elsewhere.
The reference-metric denominator in a larger domain is used
for the deterministic lower bound and need not be retained
as an unchanged output passage observation.

This use of restriction compatibility is the reason the real and
imaginary dependency patches are specified separately.
\end{remark}

\begin{proposition}[Rigidity from local forcing]
\label{comp:rigidity-implication}
Suppose the comparison family satisfies the hypotheses of
Proposition~\ref{comp:equality-tangent}.  Assume in addition that the
positive-probability rectangular shortcuts can be amplified to a
high-probability shell event with a finite deterministic catalogue
satisfying Lemma~\ref{comp:gate-forcing}.  The event must include the
following data, with constants independent of the scale:
\begin{enumerate}[label=(\alph*)]
\item regular upper and lower reference bounds of order
$w=(M\lambda^{-i})^p$, and a shortcut with denominator in
$[a'\lambda^{-ip},A'\lambda^{-ip}]$;
\item shifts of uniformly bounded amplitude and Dirichlet energy,
constant on the complete real-field dependency patch of the retained
shortcut, vanishing on the preserved exterior patch;
\item the likelihood clipping of
Lemma~\ref{comp:restricted-cm-transfer}, imposed before applying
Lemma~\ref{comp:ann-coverage};
\item output-observable shell traversal events with a lower reference
cost $\kappa w$ and the bounded overlap and alignment multiplicities
in Remark~\ref{comp:counting-checks}.
\end{enumerate}
Assume the exact local product kernels permit the completion and
preservation in Lemma~\ref{comp:restricted-cm-transfer}.
Then the optimal comparisons coincide, $c=C$, and
$F_U=CD_U$ for every open $U$.
\end{proposition}

\begin{proof}
If $c<C$, take the equality tangent and the positive-probability event
in Equation~\eqref{comp:fixed-saturation}.  Restrict its interval length
below by a deterministic positive number and its total anchor length
above by a deterministic finite number; some such restriction still
has positive probability.  Apply Lemma~\ref{comp:ann-coverage} to the
favorable finer labels given by Proposition~\ref{comp:amplification},
with volume multiplier $M$ and with the full good event (a)--(c).

On the retained middle subarc, every point lies in at least $\beta N$
good holes.  A good ball accounts for at most $K_0w$ of geodesic length:
the time span between its first and last visits is bounded by a
reference boundary connection, since a geodesic minimizes between
any two of its points.  Integration over the subarc gives
\[
 \sum_{\text{good hit balls}}w\ge K_0^{-1}\beta N
       \operatorname{len}_D(\text{middle subarc}).
\]
Choose a recipe by a fixed ordering among the successful members of
the finite catalogue.  This partitions source events and adds no
probability conditioning factor.  Intersect with the deterministic
uniform-error events of
Remark~\ref{comp:deterministic-stability-errors}; their lost probability
tends to zero.  The expected source weight is still at least
$\alpha_0N$ for some $\alpha_0>0$.

For each deterministic candidate index, perform its fixed shift and
complete the passage law before restricting to the selected source
event.  Lemma~\ref{comp:restricted-cm-transfer} gives an output
submeasure $\nu_\beta\le\Lambda\mu$ with the same fixed common law
$\mu$.  Gate forcing gives a strict shortcut on the output geodesic.
Shrinking-modification stability puts the gain in the same fixed
fractional interval and shows that the output ratio from Equation~\eqref{comp:output-ratio}
belongs to
\[
 H_N=\{C-\epsilon_N\le R<C\},\qquad \epsilon_N\downarrow0.
\]
Here $R$ is the same function of the output fields and passage data,
and $\mu$ is the common dominating law.

The output traversal tests in (d), bounded overlap at each spatial
radius, $O(N)$ radii, and boundedly many aligned labels and recipes
give total output weight at most $K_1N$ times its anchor length.
The source length cutoff and the common amplitude bound give a
deterministic output length cutoff.  Thus the hypotheses of
Lemma~\ref{comp:weighted-output-multiplicity} hold and force
$\mu(H_N)\ge\alpha_0/(\Lambda K_2)>0$.  Bounded convergence gives
$\mu(H_N)\to0$, a contradiction.

We have $c=C>0$.  The upper bound is
Equation~\eqref{comp:upper-internal}, and the lower bound follows from
Lemma~\ref{comp:intrinsic-lower}.  Therefore $F_U=CD_U$.
\end{proof}

The catalogue-availability hypothesis in
Proposition~\ref{comp:rigidity-implication} combines
Proposition~\ref{comp:conditional-catalogue-availability} with
Lemma~\ref{comp:finite-catalogue-construction}.  The real-field dependency
patch must fit inside the denominator collar: preserving a shortcut
on a much larger constant-shift patch could allow a bypass around
that collar.  An enlarged imaginary-field buffer is harmless because
the imaginary field is unchanged.  This distinction is part of the
local-kernel requirement in this application.

\begin{lemma}[Verification of the forcing hypotheses]
\label{common:forcing-verification}
Let a family satisfy Definition~\ref{common:admissible}, the reference
metric inputs above, and the finite positive comparison hypotheses of
Proposition~\ref{comp:equality-tangent}.  Then its equality tangent
satisfies the catalogue, clipping, completion, and output-count
hypotheses of Proposition~\ref{comp:rigidity-implication}.
\end{lemma}
\begin{proof}
The low event is Proposition~\ref{comp:amplification}(ii), with its
closed rectangular envelope and actual internal denominator metric.
The kernel restriction and product properties give condition (i) of
Proposition~\ref{comp:conditional-catalogue-availability}; in the
one-order case its local order event is the whole sample space.
Otherwise Definition~\ref{common:admissible}(vi) and
Lemma~\ref{comp:local-order-certificate} provide that event using a
larger imaginary buffer.  The exact volume and affine identities give
condition (iii), and Lemma~\ref{common:disk-profiles} gives condition
(iv).  The reference internal box and face bounds give condition (ii).
Thus the complete local catalogue, including fresh observations at each
exact shift and likelihood clipping, has probability as close to one
as required by Lemma~\ref{comp:ann-coverage}.

We verify explicitly how a hit of the old reference geodesic supplies
the deterministic forcing data.  Write $u=\lambda^{-ip}$ and
$w=M^p u$.  Include in the good event a lower reference crossing bound
$aw$ between the central hole and the outer circle, a surrounding
boundary connection bound $Aw$, and fixed reference connection moduli.
On a deep hit, the reference geodesic pays at least $aw$ between first
entry and last exit from the outer circle.  Its exterior pieces pay
at least the two exterior infima $\ell_0,\ell_1$.  Hence
\[
 D(a_0,a_1)\ge\ell_0+\ell_1+aw.
\]
Its exterior minimizing landing points have reference distance at
least $aw$, by concatenating their exterior approaches.  The truncated
reference modulus separates them by a fixed positive Euclidean amount,
so they fit a finite outer-cap grid with the required small errors.

Let the prepared low denominator satisfy $d\ge a'u$.  Choose
$e=\varepsilon u$ with $4C\varepsilon<\delta_0a'$.  The thin
preservation patches, terminal caps, and their connection bound $A_1u$
have already been fixed in the catalogue-availability proof.  Choose
the fixed $M$ so large that
\[
 A_1u+4e<aw,
 \qquad B=A_1u+2e.
\]
All later core, gap and access bounds concern a finite family of
new domains and affect only wall heights, core heights and smoothing
widths.  Lemma~\ref{comp:finite-catalogue-construction} therefore
applies without changing $A_1$ or the terminal errors.  Its budget and
gate conditions give $B+2e<aw$ and $4Ce<\delta_0d$, precisely the
strict inequalities of Lemma~\ref{comp:gate-forcing}.
The resulting strict shortcut has compact confinement in the image
of $B_3(0)$ around the hit ball.  These neighborhoods shrink uniformly
to the old compact middle subarc and eventually lie in its fixed
saturation neighborhood $U$.

Fix the Gaussian pinning circle outside that neighborhood and all
the late candidate supports, as in the equality-tangent proof.
Every deterministic recipe is then a Cameron--Martin direction of
this same actual Gaussian law.  Its real shift is constant on the
complete preservation patch and vanishes on the preserved exterior
patch; the larger imaginary buffer is unchanged.  Resample the local
area marks with their shifted intensity.  The conditional product
kernel and standard Borel completion in
Definition~\ref{common:admissible} give exactly the unselected
construction of Lemma~\ref{comp:restricted-cm-transfer}.  Its finite
Gaussian pairing cutoffs are part of the catalogue event before
annulus coverage or a geodesic hit is selected.

For clarity, the required output count can use a shell where the
shift is identically zero.  Include among the local reference tests
a lower bound $\kappa w$ for every crossing of
$B_{3r}(z)\setminus\overline{B_{2r}(z)}$.
Every forced reference geodesic enters the compartment inside
$B_{2r}(z)$ from anchors outside $B_{4r}(z)$, and crosses this
unchanged shell.  Its output length inside $B_{4r}(z)$ is therefore
at least $\kappa w$.  Define the output event by this length bound,
the inherited deterministic anchor-length cutoff, and the alignment
condition on the outer circle average.  That average is unchanged
by the recipe supported in $B_{2r}(z)$, so every generated output
satisfies the event.  These are functions only of the output field,
its fixed-anchor geodesic and the deterministic candidate index.

For each fixed ball, alignment allows at most
$2+\lceil2L/\log\lambda\rceil$ volume labels, and there are only
finitely many recipes.  The center mesh has bounded overlap at each
radius, and Lemma~\ref{comp:ann-coverage} uses $O(N)$ radii.
Charging the output event's weight to its geodesic length in
$B_{4r}(z)$ thus bounds the total output weight by a constant times
$N$ times the output anchor length.  The source length cutoff and
uniform shift amplitude bound this last length deterministically.
This proves (d) of Proposition~\ref{comp:rigidity-implication}; the
preceding paragraphs prove (a)--(c) and its completion hypothesis.
Uniform stability follows from
Lemma~\ref{comp:shrinking-modification-stability}.
\end{proof}

\begin{theorem}[Identification of local passage laws]
\label{common:identification}
Fix $\gamma\in[\sqrt2,2)$ and $\lambda>1$.  Let a family of local
passage laws satisfy Definition~\ref{common:admissible} and the
reference metric inputs stated above.  Suppose a law in this family
has a finite positive deterministic optimal upper comparison
$0<C<\infty$, with its comparison bounds retained at its fixed
quantum offsets and further extractions.  Then its optimal lower
comparison equals $C$, and
\[
 F_U(x,y)=C D_h(x,y;U)
\]
on every ordinary open chart $U$, with compact confinement and
exhaustion as above.  The multiplier $C$ is deterministic for the
entire joint law.  This theorem gives no estimate for a prescribed
discrete endpoint.
\end{theorem}
\begin{proof}
If the lower optimal constant $c$ were smaller than $C$, the common
annulus lemma, comparison transport, and amplification would give the
equality tangent of Proposition~\ref{comp:equality-tangent}.
Lemma~\ref{common:forcing-verification} supplies every additional
hypothesis of Proposition~\ref{comp:rigidity-implication}, whose
weighted output argument is a contradiction.  Thus $c=C>0$.
The intrinsic lower bound is Lemma~\ref{comp:intrinsic-lower}; the
upper bound follows from continuum joining and port comparison,
as in Equation~\eqref{comp:upper-internal}.  Countable compact
exhaustion gives the simultaneous assertion on open charts.
The constants in Equation~\eqref{comp:optimal-constants} are
deterministic properties of the whole law; proving them equal is
what excludes an additional sample-dependent multiplier.
\end{proof}
